\documentclass[11pt,a4paper]{article}

\usepackage[T1]{fontenc}
\usepackage[utf8]{inputenc}
\usepackage[brazilian]{babel}
\usepackage{geometry}
\usepackage{booktabs}
\usepackage{longtable}
\usepackage{xcolor}
\usepackage[hidelinks]{hyperref}

\geometry{margin=2.5cm}
\setlength{\parindent}{0pt}
\setlength{\parskip}{0.6em}

\title{LCQUI\\Estratégia e Evidências de Verificação da Implementação}
\author{Projeto LCQUI}
\date{Baseline de 26 de setembro de 2026}

\begin{document}
\maketitle

\begin{center}
\fcolorbox{black}{gray!10}{\parbox{0.9\textwidth}{
\textbf{Natureza deste documento.} Este documento não constitui fonte normativa
do comportamento do LCQUI. A autoridade normativa permanece na especificação
principal. Este documento registra estratégia, rastreabilidade e evidências de
que a implementação é verificada contra aquela especificação.
}}
\end{center}

O documento principal responde o que o LCQUI deve fazer. Este documento responde
como sabemos que a implementação satisfaz a especificação, sem redefinir o
comportamento esperado.

\tableofcontents
\clearpage

\section{Estratégia}

\subsection{Oráculo e precedência}

A verificação parte da seguinte precedência:

\begin{center}
documentação normativa $\rightarrow$ CUE $\rightarrow$ Alloy $\rightarrow$
matriz de implementação $\rightarrow$ implementação
\end{center}

Os worklogs M0--M13 registram rationale e contraexemplos, mas não substituem a
redação normativa final. A implementação existente também não altera o contrato
silenciosamente. Um conflito real entre fontes normativas interrompe o trabalho
até que seja relatado e resolvido na fonte competente.

\subsection{Camadas}

\begin{itemize}
  \item testes unitários isolam funções puras ou dependências controladas por mocks;
  \item testes de integração exercitam Admin SDK, callables e serviços locais;
  \item testes de Security Rules usam os Emulators e contextos autenticados;
  \item testes E2E exercitam fluxos observáveis no navegador;
  \item build, typecheck e lint fornecem evidência estrutural, não conformidade funcional.
\end{itemize}

\subsection{TDD incremental e RED/GREEN}

A unidade de trabalho é uma ou poucas features \texttt{IMP-*} fortemente
relacionadas. Primeiro se extrai o contrato e se cria um teste que falha pela
ausência ou divergência pretendida (RED). A menor mudança de implementação é
então feita até o teste passar (GREEN), seguida da regressão relevante. Uma
falha de ambiente, dependência ou comando é \emph{INFRA\_FAILURE}; ela não é um
RED válido de produto.

Uma feature somente pode mudar para \texttt{NÃO DIVERGENTE} depois da revisão
das obrigações normativas e da aprovação das camadas aplicáveis de backend,
Rules, frontend e E2E, sem divergência conhecida restante. Esse estado não
equivale a homologação de produção.

\subsection{Limitações}

Teste aprovado demonstra apenas os cenários cobertos. Mocks não provam integração;
compilação não prova invariantes; e uma suite legada pode codificar contrato
anterior. Resultados devem registrar versão, comando, escopo e classificação da
falha. Ausência de teste não é evidência de conformidade.

\section{Rastreabilidade}

A cadeia mínima de evidência é:

\begin{center}
Norma $\downarrow$ \texttt{IMP-*} $\downarrow$ \texttt{TEST-*}
$\downarrow$ arquivo de teste $\downarrow$ resultado
\end{center}

Cada caso permanente recebe um identificador estável por camada:

\begin{itemize}
  \item \texttt{TEST-UNIT-*}: lógica isolada;
  \item \texttt{TEST-INT-*}: integração com serviços locais;
  \item \texttt{TEST-RULES-*}: autorização em Firestore ou Storage Rules;
  \item \texttt{TEST-E2E-*}: fluxo observável no navegador.
\end{itemize}

O sufixo deve ser legível e estável. Quando houver milestone relevante, ele pode
ser incorporado, por exemplo \texttt{TEST-RULES-M9-001}. O registro associa o ID
à referência normativa exata, à feature da matriz, ao arquivo/símbolo exercitado,
ao resultado e à evidência da execução.

\subsection{Exemplo não executado}

O encadeamento abaixo é somente um exemplo de identificação futura; não declara
teste criado nem resultado:

\begin{center}
Seção 7, M9 $\rightarrow$ \texttt{IMP-RULES-001}
$\rightarrow$ \texttt{TEST-RULES-M9-001}
$\rightarrow$ futuro arquivo de Rules $\rightarrow$ resultado ainda não produzido
\end{center}

Não se preenchem resultados futuros e não se cria rastreabilidade retroativa por
inferência. A evidência deve nascer da execução registrada.

\section{Backend}

Os testes de backend são derivados em conjunto de CUE, Alloy, Seções normativas
e da divergência descrita na matriz.

\begin{longtable}{p{0.28\textwidth}p{0.62\textwidth}}
\toprule
Fonte & Tradução para cenário concreto \\
\midrule
\endhead
Fixture CUE válida & cenário positivo com payload, precondições e estado posterior esperados \\
Fixture CUE inválida & rejeição explícita, sem efeitos parciais \\
Witness Alloy SAT & cenário positivo alcançável, materializado com entidades e relações concretas \\
Assertion ou propriedade Alloy & um ou mais cenários negativos, de transição ou concorrência \\
Seção normativa e matriz & autorização, efeitos, falhas, auditoria e lacuna específica da feature \\
\bottomrule
\end{longtable}

Um resultado UNSAT não é copiado mecanicamente para Jest. A propriedade formal
é traduzida em arranjo, operação, interleaving quando aplicável e assertions
observáveis. Jest é usado para lógica/callables; Firebase Emulator Suite, para
integração; e \texttt{@firebase/rules-unit-testing}, para Rules.

\subsection{Suites existentes}

Na baseline desta rodada, \texttt{auth.test.ts} e \texttt{domain/*.test.ts} são
unitários/mocados. Os testes de integração sob \texttt{integration/} e os testes
de notificações, patrimônio, reagentes, relatórios, roteiros e turmas usam Admin
SDK ou \texttt{firebase-functions-test} em processo e dependem dos Emulators.
Nenhum deles chama o Functions Emulator por HTTP.

Os 46 casos atuais de \texttt{relatorios.test.ts} contêm assertions triviais e
são classificados como testes estruturais/textuais de baixo valor. Permanecem
na baseline, mas um PASS nesses casos não prova os comportamentos descritos.

\section{Security Rules}

Security Rules são verificadas com \texttt{@firebase/rules-unit-testing} contra
Firestore e Storage Emulators locais. Casos devem cobrir, conforme aplicável:

\begin{itemize}
  \item acesso permitido e acesso negado;
  \item o mesmo ato por outro usuário;
  \item o mesmo usuário sobre outro recurso;
  \item token ou claims obsoletos;
  \item conta inativa, versão divergente e papel persistido divergente;
  \item documento de autoridade ausente e estado alterado entre leitura e commit.
\end{itemize}

Dados de preparação são inseridos com Rules desabilitadas somente no bloco de
arranjo. A operação sob teste usa contexto autenticado ou não autenticado real.
Nenhum teste desta camada pode atingir um projeto Firebase remoto.

A primeira feature planejada para este pipeline é \texttt{IMP-RULES-001}, sobre
autoridade de usuário, versão e papéis. Ela não foi implementada nem corrigida
nesta rodada.

\section{E2E}

Os casos E2E derivam principalmente da Seção 9 (fluxos), da Seção 8 (telas e
estados) e da Seção 7 (resultados e regras). O estado atual da interface não é
fonte para inventar comportamento.

Playwright CLI é ferramenta de exploração: snapshots, locators, requests,
console, screenshots e depuração. Playwright Test mantém specs, assertions e
regressão automatizada. A exploração não substitui a suite permanente.

Locators permanentes preferem papel, label e texto semanticamente estável;
\texttt{data-testid} é usado quando necessário. Seletores CSS baseados na
estrutura visual devem ser evitados. Sessões exploratórias independentes não
compartilham estado sem intenção explícita.

O bootstrap Playwright 1.61.1 com browsers Nix está disponível, mas esta rodada
não criou configuração de suite nem teste E2E de negócio. Esses artefatos devem
nascer com o primeiro fluxo real.

\section{Resultados da baseline}

\subsection{Estados}

\begin{description}
  \item[PASS] comando executou e todas as verificações abrangidas passaram;
  \item[FAIL] comando executou e preservou uma ou mais falhas;
  \item[NOT\_RUN] não executado nesta baseline;
  \item[BLOCKED] pré-condição externa impediu a execução;
  \item[INFRA\_FAILURE] ferramenta, dependência ou ambiente impediu observar o produto;
  \item[KNOWN\_IMPLEMENTATION\_DIVERGENCE] falha coerente com divergência já registrada na matriz;
  \item[LEGACY\_TEST\_EXPECTATION] fixture ou expectativa codifica interface/estado anterior e requer rederivação normativa.
\end{description}

\subsection{Execução de 26 de setembro de 2026}

Ambiente observado: Node 26.8.1, npm 11.19.0, Java 21.0.12,
Jest 29.7.0, Firebase CLI 15.29.0 e Playwright 1.61.1. O backend declara Node 20;
o aviso de engine produzido pelo Node 26 fica registrado e não impediu a baseline.

\begin{longtable}{p{0.22\textwidth}p{0.14\textwidth}p{0.54\textwidth}}
\toprule
Verificação & Estado & Evidência \\
\midrule
\endhead
Backend build & PASS & \texttt{npm run build}; exit 0. \texttt{functions/lib} permaneceu ignorado. \\
Backend lint & FAIL & \texttt{npm run lint}; ESLint local iniciou e encontrou 82 erros existentes, sem falha de infraestrutura. \\
Jest unitário & PASS & \texttt{npm run test:unit}; 3 suites, 53 testes aprovados, 0 falhas. \\
Rules & PASS & \texttt{npm run test:rules}; 2 suites, 48 testes aprovados, 0 falhas; Emulators encerrados. \\
Integração & FAIL & \texttt{npm run test:integration}; 7 suites: 3 passaram e 4 falharam; 69 testes passaram e 16 falharam. \\
Wrapper agregado & FAIL & \texttt{npm run test:emulator}; Rules passaram, integração preservou exit 1 e os Emulators foram encerrados. \\
Frontend typecheck & PASS & \texttt{npx --no-install tsc --noEmit}; exit 0. \\
Frontend lint & FAIL & \texttt{npm run lint}; 79 erros e 29 avisos existentes. \\
Playwright bootstrap & PASS & Dependências e browser Nix 1.61.1; smoke aprovado na rodada anterior, não repetido. \\
\bottomrule
\end{longtable}

\subsection{Classificação das 16 falhas de integração}

Treze falhas são classificadas como \texttt{LEGACY\_TEST\_EXPECTATION} nesta
baseline: dez fixtures de reagentes não fornecem \texttt{idResumoReagente}; uma
expectativa patrimonial exige texto antigo de validação; e dois casos de usuário
não preparam a autoridade ativa agora exigida antes da regra que pretendem testar.
Essa classificação não autoriza atualizar os testes sem rederivá-los da norma.

Três falhas são \texttt{KNOWN\_IMPLEMENTATION\_DIVERGENCE}: uma em patrimônio e
duas em compartilhamento de roteiro executam leitura depois de escrita na mesma
transação Firestore. As áreas correspondentes já estão DIVERGENTES na matriz.
Nenhuma causa foi corrigida nesta rodada.

Não houve \texttt{INFRA\_FAILURE} após a instalação local do ESLint e a criação
dos wrappers. A baseline vermelha representa código/expectativas existentes e
continua visível pelos exit codes.

\section{Evidência de IMP-RULES-001}

\textbf{Feature:} \texttt{IMP-RULES-001} --- autoridade de usuário, versão e
papéis nas Firestore Security Rules.

\textbf{Estado anterior:} \texttt{DIVERGENTE}. Os helpers aceitavam autenticação
ou papel declarado na claim sem consultar atividade, versão e papel persistidos;
as coleções de papel também tinham leitura ampla para qualquer autenticado.

\subsection{Contrato e fontes}

Foram aplicadas a Seção 7/M9, que é a fonte normativa única de autorização, a
interpretação obrigatória de M9 e a tabela de coleções da Seção 11, a shape
\texttt{\#M9Usuario}/\texttt{\#M9Papel} do CUE e as propriedades e witnesses de
\nolinkurl{specification/alloy/operations/authorization_m9.als}. O escopo desta
evidência é a autoridade base: autenticação, existência e atividade do usuário,
versão corrente, conjunto fechado de papéis, concessão persistida e superfície
das coleções de autoridade. Escopo de almoxarifado, ACL de domínio e Storage
permanecem em \texttt{IMP-RULES-002--004}.

\subsection{Casos permanentes e ciclo RED/GREEN}

O arquivo de teste dedicado é:

\begin{center}\small
\path{functions/src/__tests__/security/m9-authority.rules.test.ts}
\end{center}

Ele usa \path{@firebase/rules-unit-testing} contra o Firestore Emulator. Dados
privilegiados são inseridos com Rules desabilitadas somente no arrange; cada
operação observada usa contexto autenticado ou não autenticado real.

\small
\begin{longtable}{@{}p{0.19\textwidth}p{0.45\textwidth}p{0.12\textwidth}p{0.11\textwidth}@{}}
\toprule
ID & Obrigação observada & Rules antigas & Rules novas \\
\midrule
\endhead
TEST-RULES-M9-001 & Não autenticado nega & PASS controle & PASS \\
TEST-RULES-M9-002 & UID sem \texttt{Usuarios} nega & FAIL esperado & PASS \\
TEST-RULES-M9-003 & Usuário inativo nega & FAIL esperado & PASS \\
TEST-RULES-M9-004 & Token sem versão nega & FAIL esperado & PASS \\
TEST-RULES-M9-005 & Versão obsoleta nega & FAIL esperado & PASS \\
TEST-RULES-M9-006 & Versão e papel persistido coerentes permitem & PASS controle & PASS \\
TEST-RULES-M9-007 & Ausência de papel nega & FAIL esperado & PASS \\
TEST-RULES-M9-008 & Papel desconhecido nega & FAIL esperado & PASS \\
TEST-RULES-M9-009 & Claim sem papel persistido nega & FAIL esperado & PASS \\
TEST-RULES-M9-010 & Papel persistido incompatível com token nega & FAIL esperado & PASS \\
TEST-RULES-M9-011 & Leitura de papel limita próprio/terceiro/query & FAIL esperado & PASS \\
TEST-RULES-M9-012 & Coleção interna permanece negada & PASS controle & PASS \\
TEST-RULES-M9-013 & Versão persistida inválida nega & FAIL esperado & PASS \\
TEST-RULES-M9-014 & \texttt{id\_usuario} do papel incompatível nega & FAIL esperado & PASS \\
\bottomrule
\end{longtable}
\normalsize

No RED funcional, 11 de 14 casos falharam pela aceitação indevida das Rules
antigas; os três restantes eram controles já satisfeitos. Após a implementação,
\texttt{TEST-RULES-M9-001--014} passou integralmente.

\subsection{Mudança e regressão}

\texttt{firestore.rules} passou a exigir documento
\nolinkurl{Usuarios/uid} existente e ativo, versão persistida inteira e não
negativa igual à claim, papel pertencente ao conjunto fechado e documento do
papel no mesmo UID com \texttt{id\_usuario} correspondente. As coleções
\texttt{Professor} e \texttt{Aluno} permitem o documento próprio ou Chefia; as
demais coleções de papel permitem somente o documento próprio. Escritas de
papel e coleções internas continuam negadas ao cliente.

As expectativas dos testes legados não foram alteradas. Seu arrange recebeu
somente fixtures M9 válidas para que cada caso continuasse medindo sua ACL de
domínio original. A regressão \texttt{npm run test:rules} passou: 3 suites,
62 testes aprovados e nenhuma falha. A regressão unitária
\texttt{npm run test:unit} passou: 3 suites e 53/53 testes.

\subsection{Ambiente e estado final}

As execuções usaram Node 26.8.1; o backend declara Node 20. A baseline de Rules
antes da mudança passou com 48/48 e o Emulator aceitou toda a campanha, logo
\texttt{NODE26\_IMPACT = NONE\_FOR\_RULES}. Isto não torna Node 26 equivalente
ao runtime de produção de Cloud Functions: evidências futuras de Functions
devem ser reproduzidas em Node 20.

\textbf{Estado final:} \texttt{NÃO DIVERGENTE} para as obrigações examinadas de
\texttt{IMP-RULES-001}. Isso não significa homologação de produção e não altera
o estado de \texttt{IMP-RULES-002}, \texttt{IMP-RULES-003} ou
\texttt{IMP-RULES-004}.

\section{Evidência de IMP-ROLE-003 --- Etapa A (fundação M9 backend)}

\textbf{Feature:} \texttt{IMP-ROLE-003} --- autorização persistida no commit.

\textbf{Estado anterior:} \texttt{DIVERGENTE}. O backend autorizava por
\texttt{validarPermissao}, que decide somente pela claim JWT; a maioria das
mutações não relia atividade, versão e papel persistidos, e o escopo de domínio
era consultado fora da unidade transacional.

\textbf{Estado depois:} \texttt{DIVERGENTE}. Esta rodada entrega apenas a
fundação de autoridade persistida M9 e migra o primeiro conjunto restrito de
operações de usuários/papéis. As demais callables ainda não foram migradas, logo
a feature transversal permanece aberta e as contagens globais da matriz não
mudam.

\subsection{Contrato M9 derivado}

A decisão de autoridade base exige, de forma fail closed: autenticação;
\texttt{Usuarios/uid} existente e \texttt{ativo=true}; \texttt{versao\_permissoes}
persistida inteira e não negativa; claim \texttt{versao\_permissoes} presente e
igual à persistida; claim \texttt{roles} em lista; toda claim de papel
pertencente ao conjunto fechado e com concessão persistida cujo
\texttt{id\_usuario} coincide com o UID; e o papel exigido pela operação presente
entre os papéis afirmados e persistidos. Fontes: Seção~7/M9, a interpretação
normativa da Seção~10.2--10.4 (\textit{erratum} M9), CUE
\texttt{formal\_m9\_autorizacao} e Alloy M9
(\texttt{M9-INV-001/002/005/008}, \texttt{M9-WIT-014/017/018}) e a superfície já
fixada em \texttt{IMP-RULES-001}. Escopo específico de domínio (vínculo de
almoxarifado, ownership de professor) permanece fora da primitiva base.

\subsection{Primitiva implementada}

\begin{center}\small
\path{functions/src/auth.ts}
\end{center}

\begin{itemize}
  \item \texttt{PAPEIS\_CONHECIDOS}: conjunto fechado M9 como fonte única no
  backend, reutilizado por \path{functions/src/usuarios.ts}.
  \item \texttt{extrairClaimsAutoridade(request)}: valida apenas a estrutura das
  claims (auth, lista de papéis, conjunto fechado, versão inteira não negativa).
  \item \texttt{resolverAutoridadePersistidaTx(tx, claims, papeisRequeridos)}:
  núcleo transacional; lê, na mesma transação do efeito, \texttt{Usuarios/uid} e
  cada documento de papel, e decide. Retorna
  \texttt{\{uid, papeis, versaoPermissoes, papelAutorizado\}}.
  \item \texttt{validarAutoridadePersistida}/
  \texttt{validarAutoridadePersistidaComClaims}: conveniências para pré-checagem
  antes de efeitos externos, reexecutando a leitura numa transação.
  \item \texttt{validarPermissao} permanece marcada como \textbf{LEGACY AUTH
  PATH}, explicitamente insuficiente para mutações críticas M9.
\end{itemize}

\subsection{Casos permanentes e ciclo RED/GREEN}

\begin{center}\small
\path{functions/src/__tests__/integration/m9-backend-authority.integration.test.ts}
\end{center}

Os casos usam Firestore Emulator e Admin SDK (não são mocks), com claims e
\texttt{request} construídos determinísticamente. A tabela registra a fonte
normativa, o comportamento do mecanismo anterior (RED) e o resultado após a
implementação (GREEN).

\small
\begin{longtable}{@{}p{0.18\textwidth}p{0.38\textwidth}p{0.13\textwidth}p{0.11\textwidth}@{}}
\toprule
ID & Obrigação observada & Antes & Depois \\
\midrule
\endhead
TEST-INT-M9-AUTH-001 & Não autenticado nega & PASS controle & PASS \\
TEST-INT-M9-AUTH-002 & UID sem \texttt{Usuarios} nega & FAIL & PASS \\
TEST-INT-M9-AUTH-003 & Usuário inativo nega & FAIL & PASS \\
TEST-INT-M9-AUTH-004 & Claim sem versão nega & PASS controle & PASS \\
TEST-INT-M9-AUTH-005 & Versão obsoleta nega & FAIL & PASS \\
TEST-INT-M9-AUTH-006 & Versão corrente + papel persistido autoriza & PASS controle & PASS \\
TEST-INT-M9-AUTH-007 & Papel permitido ausente nega & PASS controle & PASS \\
TEST-INT-M9-AUTH-008 & Papel desconhecido nega & PASS controle & PASS \\
TEST-INT-M9-AUTH-009 & Claim de papel sem concessão persistida nega & FAIL & PASS \\
TEST-INT-M9-AUTH-010 & Papel persistido não afirmado na claim nega & PASS controle & PASS \\
TEST-INT-M9-AUTH-011 & Documento de papel com UID incompatível nega & FAIL & PASS \\
TEST-INT-M9-AUTH-012 & \texttt{roles} não-array nega & PASS controle & PASS \\
TEST-INT-M9-AUTH-013 & Multi-role válido (Aluno+Bolsista) autoriza & PASS controle & PASS \\
TEST-INT-M9-AUTH-014 & Multi-role com papel desconhecido nega & PASS controle & PASS \\
\bottomrule
\end{longtable}
\normalsize

O RED foi executado antes da implementação, contra o mecanismo anterior: 6 de 19
casos falharam (002, 003, 005, 009, 011 e OP-002), evidenciando precisamente a
aceitação de token sem atividade, sem versão corrente e sem papel persistido.
Após a implementação, os 19 casos passam.

\subsection{Operação real protegida e ausência de escrita parcial}

Além da primitiva, \path{functions/src/usuarios.ts} foi migrado para a nova
autoridade transacional nas operações \texttt{convidarUsuario},
\texttt{revogarUsuarioPapel} e no núcleo \texttt{alterarPapel}. O bloco seguinte
do mesmo arquivo de teste prova a operação real.

\small
\begin{longtable}{@{}p{0.24\textwidth}p{0.58\textwidth}@{}}
\toprule
ID & Resultado \\
\midrule
\endhead
TEST-INT-M9-AUTH-OP-001 & Chefia corrente e persistida autoriza; revogação aplicada. \\
TEST-INT-M9-AUTH-OP-002 & Token obsoleto nega e o Firestore permanece inalterado. \\
TEST-INT-M9-AUTH-OP-003 & Chefia inativa nega e o Firestore permanece inalterado. \\
TEST-INT-M9-AUTH-OP-004 & Claim de Chefe sem documento persistido nega e não muta. \\
\bottomrule
\end{longtable}
\normalsize

Cada caso negado verifica o estado persistido posterior, não apenas o
\texttt{HttpsError}: o documento do papel do alvo e o \texttt{ativo} permanecem.

\subsection{Ambiente e estratégia Node 20}

O backend declara \texttt{engines.node=20}. A evidência final foi executada em
Node \textbf{20.19.6} (npm 10.8.2). O \texttt{nixpkgs} atual removeu
\texttt{nodejs\_20} por EOL em 30/04/2026; a execução usou
\path{nix shell github:nixos/nixpkgs/nixos-25.05#nodejs_20} como fornecimento
temporário, sem alterar \texttt{flake.lock} e sem alterar o Home Manager:

\begin{center}
\texttt{NODE20\_STRATEGY = TEMPORARY\_NIX\_ONLY}; \quad
\texttt{HOME\_MANAGER\_CHANGED = NO}
\end{center}

\subsection{Resultados e regressão}

\small
\begin{longtable}{@{}p{0.52\textwidth}p{0.38\textwidth}@{}}
\toprule
Verificação (Node 20.19.6) & Resultado \\
\midrule
\endhead
\texttt{npm run build} & PASS \\
\texttt{npm run test:unit} & PASS (53/53) \\
\texttt{TEST-INT-M9-AUTH-*} + \texttt{-OP-*} (emulator) & PASS (19/19) \\
\texttt{usuarios.integration.test.ts} (fixtures M9) & PASS (2/2) \\
\texttt{npm run test:integration} & 90 passam / 14 falham \\
\bottomrule
\end{longtable}
\normalsize

A baseline registrada era 69 passam / 16 falham. A diferença é integralmente
explicada: os 19 casos M9 são novos e passam, e 2 casos legados de
\texttt{usuarios.integration.test.ts} que já falhavam por fixture sem autoridade
persistida passaram a ter arrange M9 válido. As 14 falhas remanescentes estão
somente em \texttt{reagentes} (10), \texttt{patrimonio} (2) e \texttt{roteiros}
(2), nenhuma em \texttt{auth}/\texttt{usuarios}; nenhuma falha nova não explicada.

O lint direcionado aos arquivos alterados não introduz novos erros; os avisos
remanescentes (\texttt{no-explicit-any} em \texttt{catch} legado e no teste
legado) já existiam.

\subsection{Escopo e pendências}

Migrado nesta etapa: \texttt{alterarPapel}, \texttt{convidarUsuario} e
\texttt{revogarUsuarioPapel}, com a decisão de chefia relida na transação do
efeito (fecha TOCTOU). A ordem de \texttt{convidarUsuario} mantém efeitos de
Firebase Auth após uma pré-checagem M9 e antes da revalidação transacional; o
risco residual (criação de identidade Auth entre pré-checagem e commit) fica
registrado como dependência futura, sem redesenho do workflow de convite nesta
rodada.

Não migrado: reagentes, patrimônio, turmas, posts, roteiros, notificações,
relatórios, etiquetas; e \texttt{IMP-ROLE-001}, \texttt{IMP-ROLE-002} e
\texttt{IMP-IDEM-001} permanecem intocados. \texttt{IMP-ROLE-003} continua
\texttt{DIVERGENTE} por desenho.

\section{Evidência de IMP-ROLE-003 --- Etapa A.1 (endurecimento da fundação)}

\textbf{Feature:} \texttt{IMP-ROLE-003} --- autorização persistida no commit.

\textbf{Estado anterior:} \texttt{DIVERGENTE}. A fundação M9 da Etapa A relia
apenas os documentos de papel afirmados na claim, deixando invisível um papel
persistido adicional e incompatível.

\textbf{Estado depois:} \texttt{DIVERGENTE}. Esta microcorreção endurece a
primitiva para ler e validar o conjunto persistido completo. Nenhuma nova
callable foi migrada e as contagens da matriz não mudam.

\subsection{Lacuna encontrada}

A Seção~7/M9 define que papéis atuais são os documentos de papel cujo
\texttt{id\_usuario} coincide com o UID, com matriz fechada: Chefe\_Geral é
exclusivo; Professor e Aluno não coexistem; Bolsista implica Aluno; Bolsista e
Gestor\_Almoxarifado não coexistem; e estado inconsistente significa negação. A
primitiva anterior lia somente \texttt{claims.papeis}, de modo que um estado
persistido como \texttt{Chefe\_Geral + Professor} podia autorizar uma ação de
Chefia com \texttt{roles=["Chefe\_Geral"]}.

\subsection{Correção}

\path{functions/src/auth.ts}, em
\texttt{resolverAutoridadePersistidaTx}, agora:

\begin{itemize}
  \item lê os seis documentos canônicos de papel do UID a partir de
  \texttt{PAPEIS\_CONHECIDOS} (fonte única), não apenas os da claim;
  \item exige \texttt{id\_usuario == uid} em todo documento existente;
  \item executa \texttt{validarMatrizPapeis} sobre o conjunto
  \textbf{persistido} e converte combinação inválida em \texttt{permission-denied};
  \item exige que a claim projete \textbf{exatamente} o conjunto persistido na
  versão corrente (nem subconjunto nem superconjunto);
  \item só então valida o papel requerido.
\end{itemize}

\subsection{Casos permanentes e ciclo RED/GREEN}

Novos casos em
\path{functions/src/__tests__/integration/m9-backend-authority.integration.test.ts}:

\small
\begin{longtable}{@{}p{0.18\textwidth}p{0.45\textwidth}p{0.1\textwidth}p{0.1\textwidth}@{}}
\toprule
ID & Obrigação observada & Antes & Depois \\
\midrule
\endhead
TEST-INT-M9-AUTH-016 & Chefe+Professor persistidos negam (claim só de Chefe) & FAIL & PASS \\
TEST-INT-M9-AUTH-017 & Professor+Aluno persistidos negam & FAIL & PASS \\
TEST-INT-M9-AUTH-018 & Bolsista persistido sem Aluno nega & FAIL & PASS \\
TEST-INT-M9-AUTH-019 & Bolsista+Gestor\_Almoxarifado persistidos negam & FAIL & PASS \\
TEST-INT-M9-AUTH-020 & Claim incompleta (subconjunto do persistido) nega & FAIL & PASS \\
TEST-INT-M9-AUTH-021 & Claim superconjunto (papel sem persistência) nega & PASS controle & PASS \\
\bottomrule
\end{longtable}
\normalsize

RED executado contra a fundação da Etapa A: 5 de 6 casos falharam (016--020) e
021 já negava por documento ausente. Após a correção, todos passam; a campanha
M9 completa fica em 25/25.

\subsection{Node 20 e regressão}

Evidência final sob Node \textbf{20.19.6} via
\path{nix shell github:nixos/nixpkgs/nixos-25.05#nodejs_20}
(\texttt{NODE20\_STRATEGY = TEMPORARY\_NIX\_ONLY}; Home Manager inalterado).

\small
\begin{longtable}{@{}p{0.52\textwidth}p{0.38\textwidth}@{}}
\toprule
Verificação (Node 20.19.6) & Resultado \\
\midrule
\endhead
\texttt{npm run build} & PASS \\
\texttt{npm run test:unit} & PASS (53/53) \\
\texttt{TEST-INT-M9-AUTH-*} + \texttt{-OP-*} (emulator) & PASS (25/25) \\
\texttt{usuarios.integration.test.ts} & PASS (2/2) \\
\texttt{npm run test:integration} & 96 passam / 14 falham \\
\bottomrule
\end{longtable}
\normalsize

A baseline imediatamente anterior era 90/14; o delta são exatamente os 6 casos
M9 novos aprovados. As 14 falhas permanecem em \texttt{reagentes} (10),
\texttt{patrimonio} (2) e \texttt{roteiros} (2); nenhuma falha nova não
explicada.

\subsection{Pendência externa (não resolvida)}

\texttt{convidarUsuario} ainda executa efeito externo de Firebase Auth
(\texttt{createUser}) entre a pré-checagem M9 e o commit Firestore. Se a
autoridade mudar nesse intervalo, o commit Firestore é revalidado e negado, mas
uma identidade Auth pode já existir. Essa é uma divergência conhecida de
\texttt{IMP-ROLE-003}, que toca M7/outbox e \texttt{IMP-ROLE-001/002}, e
permanece como fatia futura; não foi silenciosamente absorvida nesta
microcorreção.

\textbf{Estado final:} \texttt{IMP-ROLE-003} continua \texttt{DIVERGENTE}. A
fundação agora é confiável para crescer: valida conjunto persistido completo,
matriz normativa e coerência claim/persistência.

\section{Evidência de IMP-IDEM-001 --- Etapa A (fundação canônica M7)}

\textbf{Feature:} \texttt{IMP-IDEM-001} --- receipt e retry universal de mutações
críticas.

\textbf{Estado anterior:} \texttt{DIVERGENTE}. Só \texttt{usuarios.ts} possuía
uma forma parcial de receipt, com hash \texttt{JSON.stringify(dados)}, sem
\texttt{tipo\_operacao}/\texttt{payload\_hash} canônicos, e \texttt{idOperacao}
opcional com geração de UUID novo.

\textbf{Estado depois:} \texttt{DIVERGENTE}. Esta rodada cria e comprova a
fundação canônica M7 reutilizável; nenhuma mutação de produto foi migrada e as
contagens da matriz não mudam.

\subsection{Contrato M7 derivado}

Fontes: Seção~7 (\texttt{sec:idempotencia-m7}), Seção~5
(\textit{Operações técnicas}) e o CUE \texttt{formal\_m7} (estados fechados,
\texttt{payload\_hash} hexadecimal de 64). Pontos confirmados:

\begin{itemize}
  \item identidade semântica \texttt{(uid, tipo\_operacao, payload\_hash)}; as
  três componentes são comparadas no retry;
  \item \texttt{hashPayload(tipoOperacao, payload) =
  SHA-256(tipoOperacao + "\textbackslash n" + canonicalize(payload))};
  \item \texttt{idOperacao} opaco, \texttt{[A-Za-z0-9\_-]\{1,128\}}, fora do
  hash; novo id só em nova intenção;
  \item \texttt{Operacoes/\{idOperacao\}}: \texttt{uid}, \texttt{tipo\_operacao},
  \texttt{payload\_hash}, \texttt{status} (PENDENTE/CONCLUIDA/FALHOU),
  \texttt{resultado}, \texttt{criado\_em}, \texttt{atualizado\_em};
  \item comando atômico grava \texttt{CONCLUIDA} na mesma transação; fluxo
  externo usa \texttt{PENDENTE}$\rightarrow$\texttt{CONCLUIDA}/\texttt{FALHOU};
  \item receipt corrompido falha fechado; retry não atualiza \texttt{timestamp}.
\end{itemize}

\subsection{Canonicalização}

A canonicalização ordena chaves de objeto por code unit recursivamente, ignora
\texttt{undefined} em objeto (equivalente a ausente), preserva a ordem de
arrays, distingue \texttt{null} de ausente e não aplica trim/caixa/arredondamento.
Tipos não suportados (não finitos, \texttt{bigint}, função, símbolo,
\texttt{undefined} isolado e objetos não planos como \texttt{Date}) falham
fechado, pois não pertencem ao domínio de payload JSON.

\subsection{Fundação implementada}

\begin{center}\small
\path{functions/src/idempotencia.ts}
\end{center}

\begin{itemize}
  \item \texttt{canonicalize(value)} --- impressão determinística do payload;
  \item \texttt{hashPayload(tipoOperacao, payload)} --- SHA-256 canônico global;
  \item \texttt{construirIdentidade(uid, tipoOperacao, payload)} --- monta a
  identidade semântica;
  \item \texttt{validarIdOperacao(idOperacao)} --- formato normativo;
  \item \texttt{resolverOperacaoTx(tx, idOperacao, identidade)} --- decide
  \texttt{NOVA}/\texttt{REPLAY}/\texttt{PENDENTE}/\texttt{FALHOU} ou rejeita
  reutilização incompatível com \texttt{already-exists}; receipt inconsistente
  falha fechado;
  \item \texttt{registrarOperacaoConcluidaTx}/\texttt{registrarOperacaoPendenteTx}
  --- escrita do receipt;
  \item \texttt{concluirOperacaoPendenteTx}/\texttt{falharOperacaoPendenteTx} ---
  transição \texttt{PENDENTE}$\rightarrow$\texttt{CONCLUIDA}/\texttt{FALHOU}.
\end{itemize}

\subsection{TDD: RED e GREEN}

RED contra o scaffold ingênuo (ordem de chave via \texttt{JSON.stringify},
comparação apenas de \texttt{uid}, sem validação de receipt):

\small
\begin{longtable}{@{}p{0.26\textwidth}p{0.55\textwidth}@{}}
\toprule
ID & Resultado \\
\midrule
\endhead
TEST-UNIT-M7-CANON-001 & FAIL (ordem de chave alterava o canonical) \\
TEST-UNIT-M7-CANON-002 & FAIL (objetos aninhados não ordenados) \\
TEST-UNIT-M7-CANON-009 & FAIL (tipos não suportados não falhavam fechado) \\
TEST-UNIT-M7-HASH-001 & FAIL (payload permutado gerava hash diferente) \\
TEST-INT-M7-RECEIPT-003 & FAIL (outro ator tratado como replay) \\
TEST-INT-M7-RECEIPT-004 & FAIL (outro tipo tratado como replay) \\
TEST-INT-M7-RECEIPT-005 & FAIL (outro payload tratado como replay) \\
TEST-INT-M7-RECEIPT-008 & FAIL (receipt corrompido não falhava fechado) \\
TEST-INT-M7-RECEIPT-009 & FAIL (\texttt{PENDENTE} tratado como nova) \\
TEST-INT-M7-RECEIPT-010 & FAIL (processador não implementado) \\
\bottomrule
\end{longtable}
\normalsize

Após a implementação, sob Node 20.19.6: unidade \texttt{TEST-UNIT-M7-*}
(18/18) e integração \texttt{TEST-INT-M7-RECEIPT-001--012} (12/12) aprovadas,
incluindo concorrência idêntica (\texttt{006}), intenções distintas
(\texttt{007}), \texttt{idOperacao} fora do hash (\texttt{011}) e
\texttt{FALHOU} terminal (\texttt{012}).

\subsection{Node 20 e regressão}

Evidência final sob Node \textbf{20.19.6} via
\path{nix shell github:nixos/nixpkgs/nixos-25.05#nodejs_20}
(\texttt{NODE20\_STRATEGY = TEMPORARY\_NIX\_ONLY}; Home Manager inalterado).

\small
\begin{longtable}{@{}p{0.52\textwidth}p{0.38\textwidth}@{}}
\toprule
Verificação (Node 20.19.6) & Resultado \\
\midrule
\endhead
\texttt{npm run build} & PASS \\
\texttt{npm run test:unit} & PASS (71/71, inclui 18 M7) \\
\texttt{TEST-INT-M7-RECEIPT-001--012} (emulator) & PASS (12/12) \\
\texttt{npm run test:integration} & 108 passam / 14 falham \\
\bottomrule
\end{longtable}
\normalsize

A baseline imediatamente anterior era 96/14; o delta são exatamente os 12 casos
M7 novos. As 14 falhas permanecem em \texttt{reagentes} (10),
\texttt{patrimonio} (2) e \texttt{roteiros} (2); nenhuma falha nova.

\subsection{Divergência registrada e limites}

\texttt{usuarios.ts} continua intocado e é evidência da divergência: usa
\texttt{Operacoes} com id derivado e \texttt{hash\_entrada =
JSON.stringify(dados)}, sem \texttt{tipo\_operacao}/\texttt{payload\_hash}
canônicos, e \texttt{idOperacao} opcional. Nenhum domínio de produto foi migrado;
nenhum processador/outbox externo foi implementado (apenas a taxonomia de
estados é suportada). A pendência externa de \texttt{convidarUsuario}
(Auth entre pré-checagem e commit), documentada na Etapa A.1, permanece
inalterada e é candidata natural à fundação M7/outbox posterior.

\textbf{Estado final:} \texttt{IMP-IDEM-001} continua \texttt{DIVERGENTE}. A
fundação M7 está estável e revisável isoladamente para receber a primeira
mutação real em rodada própria.


\end{document}
